\documentclass[a4paper,12pt,oneside]{article}
\usepackage{discrete-practice}

% Источник условий: https://discrete-math-practics.yonote.ru/share/663bca14-748d-4b4c-9308-ef47eff949d0/doc/dz3-konechnye-avtomaty-srok-do-14042026-2355-L5mp5WMi3m
\begin{document}
\HomeworkHeader{2}{3}{Конечные автоматы}
\noindent Срок сдачи: 14.04.2026, 23:55.\hfill
\PracticeLink{Условия в Yonote}{https://discrete-math-practics.yonote.ru/share/663bca14-748d-4b4c-9308-ef47eff949d0/doc/dz3-konechnye-avtomaty-srok-do-14042026-2355-L5mp5WMi3m}

\medskip
В ИТМО начали свою работу сервисные роботы:
\begin{numlist}
  \item Робот-бариста.
  \item Робот-гид.
\end{numlist}
У каждого робота есть заданное поведение.

\Task{Для каждого из двух роботов необходимо}
\begin{numlist}
  \item Построить \textbf{академическое} регулярное выражение, задающее этот язык.
  \item Построить $\varepsilon$-НКА, распознающий этот язык.
  \item Преобразовать $\varepsilon$-НКА в НКА.
  \item Построить эквивалентный ДКА.
  \item Записать таблицу переходов ДКА.
  \item Минимизировать ДКА.
  \item Привести примеры 3 допустимых и 3 недопустимых слов.
  \item Для одного допустимого и одного недопустимого слова показать работу автомата по шагам.
\end{numlist}

\Task{Поведение роботов: 1. Робот-бариста}
Робот принимает заказы и готовит напитки по упрощённому рецепту.

\noindent Алфавит: $\Sigma=\{o,e,m,s,p\}$, где
\begin{itemize}
  \item $o$ -- принят заказ;
  \item $e$ -- приготовлен эспрессо;
  \item $m$ -- добавлено молоко;
  \item $s$ -- добавлен сироп;
  \item $p$ -- напиток выдан.
\end{itemize}

\Task{Правила корректной работы}
\begin{numlist}
  \item Каждый напиток начинается с символа $o$.
  \item После $o$ обязательно идёт $e$.
  \item После $e$ робот может:
    \begin{itemize}
      \item сразу выдать напиток;
      \item добавить только молоко;
      \item добавить только сироп;
      \item добавить молоко, затем сироп.
    \end{itemize}
  \item После выдачи напитка $p$ может начаться следующий заказ.
  \item Слово описывает последовательность из нуля или более полностью завершенных заказов.
\end{numlist}

\noindent Пример допустимых слов: $\varepsilon$, $oep$, $oemp$, $oesp$, $oemsp$, $oepoemp$.

\noindent Пример недопустимых слов: $op$, $omp$, $eo$, $oespm$, $oemps$.

\newpage
\Task{Поведение роботов: 2. Робот-гид}
Робот встречает гостей и проводит их по маршруту.

\noindent Алфавит: $\Sigma=\{g,q,f,l,r,h\}$, где
\begin{itemize}
  \item $g$ -- робот поприветствовал гостя;
  \item $q$ -- проверил QR-пропуск;
  \item $f$ -- сделал шаг вперёд;
  \item $l$ -- повернул налево;
  \item $r$ -- повернул направо;
  \item $h$ -- маршрут завершён.
\end{itemize}

\Task{Правила корректной работы}
\begin{numlist}
  \item Каждый маршрут начинается с последовательности $gq$.
  \item После проверки пропуска робот должен совершить хотя бы одно движение.
  \item После каждого поворота $l$ или $r$ должен быть хотя бы один шаг вперед $f$.
  \item Два поворота подряд запрещены.
  \item Маршрут заканчивается символом $h$.
  \item Слово описывает \textbf{один} полностью завершенный маршрут.
\end{numlist}

\noindent Пример допустимых слов: $gqfh$, $gqffh$, $gqlfh$, $gqrffh$, $gqfflffrfh$.

\noindent Пример недопустимых слов: $gqh$, $gqlh$, $gqlrfh$, $qgfh$, $fh$.
\end{document}
